\documentclass[11pt]{article}

\usepackage[margin=1in]{geometry}
\usepackage{amsmath,amssymb,amsthm}
\usepackage{graphicx}
\usepackage{booktabs}
\usepackage{hyperref}
\usepackage{url}
\usepackage{xcolor}
\usepackage{cite}
\usepackage{authblk}
\usepackage{algorithm}
\usepackage{algpseudocode}
\usepackage{subcaption}

% Placeholder macro for experiment numbers still to be filled in.
\newcommand{\todo}[1]{\textcolor{red}{\textbf{[TBD: #1]}}}
\newcommand{\num}{--} % placeholder cell marker for tables

\newtheorem{definition}{Definition}
\newtheorem{proposition}{Proposition}
\newtheorem{remark}{Remark}

\title{Resolving Visual Observation Aliasing in\\
Vision-Language-Action Models:\\
A Feature-Level Short-Term Memory Adapter}

\author[1]{Anonymous Authors}
\affil[1]{Anonymous Institution}
\date{\today}

\begin{document}
\maketitle

\begin{abstract}
Vision-Language-Action (VLA) models for robotic manipulation are almost universally
trained as \emph{Markov} policies that predict each action from the current-frame
observation alone. We identify a specific, widespread, and silent failure mode of
this design, which we call \textbf{visual observation aliasing}: pairs of trajectory
states whose single-frame observations are nearly indistinguishable but whose expert
actions are not, forcing a Markov policy to produce mode-averaged outputs. We
formalize observation aliasing in the imitation-learning POMDP setting, propose a
data-driven \emph{aliasing score} that flags which tasks and which trajectory states
suffer from it, and show how the design of an effective memory module can be read
off directly from the structure of the problem.

Guided by this analysis we introduce the \textbf{Short-Term Memory Bank (STMB)}, a
feature-level adapter that samples only $M{=}5$ past frames at a fixed interval
$I{=}10$ and writes them into the DeepStack visual pathway of a Qwen3-VL backbone
through a per-token, per-level gated slot-attention update. We further catalogue four
previously unreported \emph{silent} failure modes of memory-augmented VLA pipelines
(training--evaluation frame mismatch, zero-image padding, write-query self-reference,
and causal confusion by history) and show that all four must be addressed for memory
to pay off. On LIBERO (\textsc{spatial}, \textsc{object}, \textsc{goal}, \textsc{10},
\textsc{90}) STMB yields \todo{overall gain}; an aliasing-stratified analysis
confirms that gains correlate with the proposed aliasing score
(\todo{correlation coefficient}), so the improvements are predictable from a
property of the data rather than from global leaderboard effects.
\end{abstract}

%=========================================================================
\section{Introduction}
\label{sec:intro}

\paragraph{The Markov assumption in VLAs is load-bearing, and brittle.}
Modern Vision-Language-Action models~\cite{brohan2023rt2,kim2024openvla,black2024pi0}
share a common simplification: the action at time $t$ is predicted from
\emph{only} the current observation $o_t$ and the language instruction~$\ell$.
This treats manipulation as a Markov Decision Process whose state is fully captured
by a single RGB frame (plus, in some variants, proprioception). The simplification
is extremely useful---it lets any pretrained VLM be repurposed as a policy with a
light action head---but it is also load-bearing in a way that few papers discuss
explicitly: once the current frame stops being a sufficient statistic of the
world state, the training loss silently pushes the policy toward the mean of the
conditional action distribution, and rollouts fail for reasons that look like
stochastic noise rather than a structural defect.

\paragraph{The problem this paper targets: visual observation aliasing.}
We give this failure mode a name. \emph{Visual observation aliasing} occurs when
two trajectory states $s_1 \neq s_2$ in the training distribution produce nearly
identical current-frame observations $o_1 \approx o_2$ but are associated with
expert actions $a_1, a_2$ drawn from meaningfully different modes
(opposite directions, opposite gripper commands, different subtask phases). Under
behaviour cloning, a Markov policy cannot distinguish $o_1$ from $o_2$ and is
forced to regress to the conditional mean $\mathbb{E}[a \mid o]$, which is often
outside the set of task-valid actions. In a pick-and-place episode the gripper
looks \emph{identical} during the reach and the retract; the robot pose looks
\emph{identical} mid-push and mid-pull; the wrist camera shows a similar patch
right before contact and right after it. These are not edge cases: in
\S\ref{sec:exp-alias} we show that in LIBERO a significant fraction of training
states participate in at least one aliasing pair.

\paragraph{Existing memory solutions either over-solve or under-solve the problem.}
A natural response is to give the policy access to more than the current frame.
The recent VLA literature takes several such steps: MemoryVLA~\cite{shi2025memoryvla}
and ReMem-VLA~\cite{remem2025} add dual perceptual--cognitive memory or recurrent
queries inside the language model; PAM~\cite{pam2025} stacks a dense 300-frame
history; LoLA~\cite{lola2025} downsamples long-horizon observations; CoT-VLA
introduces chain-of-thought visual tokens~\cite{zhao2025cotvla}. Each of these is
a substantial architectural rewrite and each targets primarily
\emph{long-horizon} or \emph{task-level} context. At the same time, a classical
result in imitation learning---\emph{causal confusion}~\cite{dehaan2019causal}---
warns that na\"ively adding observation history can \emph{hurt} performance by
introducing nuisance correlates from past actions. The existing VLA memory
literature does not reconcile these two pressures: it tends to over-parameterise
the memory (which risks causal confusion) without first localising exactly what
part of the single-frame policy is failing.

\paragraph{A problem-driven design.}
We argue that the right unit of analysis is not ``how much history can we fit''
but ``how little history is strictly required to break the aliasing, and where in
the network should it enter?'' Reading the problem back off of its definition, an
adequate fix must: (P1) intervene at the \emph{perception} level, because aliasing
is a perception-level rather than reasoning-level failure; (P2) deliver the
intervention as a \emph{feature-level} correction of the current-frame
representation, rather than as additional time-stamped tokens; (P3) use a
\emph{sparse} past window whose time scale matches the task's manipulation phases,
so that the memory provides discriminative low-frequency signal without
re-introducing causal-confusion vulnerabilities; (P4) \emph{gate} the memory
contribution per token, so that disambiguation is local (near the gripper, near
the target) rather than a global feature replacement; (P5) expose the absolute
time step, so that repetitions of the same visual motif in different task phases
can still be told apart.

\paragraph{STMB: a feature-level short-term memory adapter.}
Our proposed module, the \textbf{Short-Term Memory Bank} (STMB), is an
architecture derived from P1--P5. It maintains a tensor of past visual features at
three DeepStack levels and two camera views, runs a per-level slot-attention read
using the most recent past frame as the query, fuses the aggregated history with
the current-frame features through a learned per-token sigmoid gate, and adds a
sinusoidal encoding of the absolute time step. The fused features are written
back into the DeepStack pathway of a Qwen3-VL backbone. The entire integration
is a single conditional block in the vision--language forward pass.

\paragraph{Four silent failure modes.}
While instrumenting the memory-augmented pipeline we discovered four distinct
failure modes that silently flatten the benefits of memory. (F1)
\emph{Training--evaluation frame-index mismatch}: the dataloader collected past
frames in a convention that differed from the online rollout buffer by one
interval, so the memory seen at inference was not the memory seen at training.
(F2) \emph{Zero-image padding} at the start of an episode, which injects an
out-of-distribution black image into the memory bank where training would have
repeated the earliest real frame. (F3) \emph{Write-query self-reference}: if the
current frame is (incorrectly) included in the memory bank, the slot-attention
read uses the current frame as both query and key, which collapses the update to
a near-identity. (F4) \emph{Causal confusion by proxy}, the classical
imitation-learning pitfall that memory-augmented VLA papers rarely acknowledge.
These failures are silent because shapes, dtypes, and losses all appear correct.
We report each, show how STMB's design choices and our alignment patch defuse
them, and provide a checklist that other memory-augmented VLAs can apply.

\paragraph{Contributions.} We make the following contributions:
\begin{enumerate}
    \item \textbf{Problem formulation.} We formalize \emph{visual observation
    aliasing} in the VLA imitation-learning POMDP
    (\S\ref{sec:problem}), and we introduce a data-driven \emph{aliasing score}
    (Def.~\ref{def:alias}) that can be computed on any VLA dataset without
    training.
    \item \textbf{A problem-derived memory module.} We derive five design
    constraints P1--P5 directly from the aliasing formulation and instantiate
    them in \textbf{STMB}, a feature-level, per-token gated memory adapter over
    DeepStack features (\S\ref{sec:method}).
    \item \textbf{A diagnostic catalogue of silent failure modes} (F1--F4,
    \S\ref{sec:diagnostics}) of memory-augmented VLA pipelines, with a minimal
    code-level patch for each.
    \item \textbf{Aliasing-stratified evaluation.} Beyond reporting overall
    success rates on LIBERO, we stratify tasks by the aliasing score and show
    that gains from STMB are predicted by this property of the data rather
    than by a global leaderboard effect (\S\ref{sec:exp-alias}).
    \item \textbf{Reproducible training and evaluation pipeline} built on
    RynnBrain / Qwen3-VL with a single-file memory integration, released with
    checkpoints and configurations.
\end{enumerate}

\paragraph{Paper organization.} \S\ref{sec:related} positions our approach
against existing memory-augmented VLAs and the causal-confusion literature.
\S\ref{sec:problem} formalizes visual observation aliasing. \S\ref{sec:method}
introduces STMB and details its integration with the VLM.
\S\ref{sec:diagnostics} catalogues the four silent failure modes and documents
our alignment patch. \S\ref{sec:exp} reports LIBERO results, the aliasing-
stratified analysis, and ablations. \S\ref{sec:discussion} discusses limitations
and extensions.

%=========================================================================
\section{Related Work}
\label{sec:related}

\paragraph{Vision-Language-Action models.}
OpenVLA~\cite{kim2024openvla}, RT-2~\cite{brohan2023rt2}, $\pi_0$~\cite{black2024pi0}
and concurrent systems adapt a pretrained VLM to manipulation by attaching an
action head (discrete tokens, MLP regression, or a flow-matching expert). All of
these operate under the Markov assumption: the action at step $t$ is a function
of $(o_t, \ell)$. Our work targets exactly the regime where this assumption
breaks down.

\paragraph{Memory-augmented VLA.}
MemoryVLA~\cite{shi2025memoryvla} introduces a perceptual--cognitive dual memory
bank and a diffusion action head; ReMem-VLA~\cite{remem2025} adds frame-level and
chunk-level recurrent queries that propagate through the language model;
PAM~\cite{pam2025} addresses state ambiguity with a 300-frame dense history and
an adaptive working-memory recoding step; LoLA~\cite{lola2025} combines current
observation encoding with downsampled historical motion encoding for long-horizon
tasks. These methods are primarily designed for long-horizon or task-level
context. We instead target the \emph{perception} level where the aliasing
problem manifests, and show in \S\ref{sec:exp} that a feature-level adapter at
this layer is sufficient for the aliasing regime.

\paragraph{Causal confusion in imitation learning.}
De Haan et al.~\cite{dehaan2019causal} show that na\"ively feeding a history of
observations to a behaviour-cloning policy can \emph{hurt} performance because
the history contains nuisance correlates with past actions. Recent work
extends this analysis to robotic imitation~\cite{causalbot2025}. Our sparse
sampling ($M{=}5$, $I{=}10$) and per-token gating are inductive-bias choices
designed to admit only the discriminative low-frequency signal required by the
aliasing formulation while remaining inhospitable to the kind of fine-grained
past-action traces that drive causal confusion.

\paragraph{Visual reasoning chains for VLA.}
CoT-VLA~\cite{zhao2025cotvla} integrates chain-of-thought visual tokens into
the prompt. This adds reasoning-level context but operates on the prompt and
therefore inflates sequence length. Our approach is complementary: STMB augments
perception at feature level and can be stacked under CoT-style prompts.

%=========================================================================
\section{Problem Formulation: Visual Observation Aliasing}
\label{sec:problem}

We consider a language-conditioned manipulation task as a POMDP
$(\mathcal{S}, \mathcal{A}, \mathcal{O}, T, \Omega, \ell)$, where $\mathcal{S}$ is
the (latent) world state, $\mathcal{A}$ is the robot action space, $\mathcal{O}$
is the image observation space, $T$ is the transition kernel, $\Omega$ is the
observation kernel, and $\ell$ is a language instruction. Expert demonstrations
induce a joint distribution $p(s, o, a \mid \ell)$.

A \emph{Markov VLA policy} $\pi_\theta(a \mid o, \ell)$ is trained by imitation
against a data distribution $\widehat{p}(o, a \mid \ell)$ and in particular
makes no use of trajectory history.

\begin{definition}[Observation aliasing pair]
\label{def:alias}
Let $d(\cdot,\cdot)$ be a metric on the VLM feature space. Given a threshold
$\varepsilon > 0$ and an action-difference threshold $\delta > 0$, a pair of
training points $(o_i, a_i, \ell)$ and $(o_j, a_j, \ell)$ is an
\emph{$(\varepsilon, \delta)$-aliasing pair} if
\[
d(\phi(o_i), \phi(o_j)) \le \varepsilon \quad \text{and} \quad
\|a_i - a_j\|_2 \ge \delta,
\]
where $\phi$ is a fixed visual encoder (e.g.~the frozen Qwen3-VL DeepStack
output). The \emph{aliasing score} of an observation $o_i$ is the density of
aliasing partners in its $k$-nearest-neighbour set:
\[
\mathrm{Alias}_k(o_i) = \frac{1}{k}\sum_{j \in \mathrm{kNN}_k(o_i)}
\mathbf{1}\!\left[\|a_i - a_j\|_2 \ge \delta\right].
\]
\end{definition}

\begin{proposition}[Mode averaging under aliasing]
\label{prop:modeavg}
Let $R(\theta) = \mathbb{E}_{(o,a)}[\|\pi_\theta(o) - a\|_2^2]$ be the $L_2$
imitation objective. For any $o$, the optimal Markov policy satisfies
$\pi^*(o) = \mathbb{E}_{a \sim p(a \mid o)}[a]$. If the set
$\mathrm{Alias}_k(o)$ is high and the conditional action distribution is
multi-modal, $\pi^*(o)$ lies outside the set of task-valid actions.
\end{proposition}

This proposition makes precise the ``silent failure'' we target: at high-
aliasing observations the Bayes-optimal Markov policy is provably bad. The rest
of this paper derives the minimum context required to make the policy
non-Markov specifically at these points.

\begin{remark}[What kind of memory is enough?]
The aliasing pairs that occur in manipulation are almost always disambiguated
by \emph{low-frequency motion cues} (is the arm extending or retracting?), by
\emph{recent contact events} (did the gripper just close?), or by \emph{task
phase} (are we in the reach, transport, or release segment?). All three are
signals at the time scale of whole sub-actions, i.e.~seconds, not of control
steps. A memory module with sparse sampling at a task-relevant interval is
both \emph{necessary} (single-frame is insufficient) and \emph{sufficient}
(dense kinematic detail is not required) for this problem class.
\end{remark}

%=========================================================================
\section{Method: Short-Term Memory Bank (STMB)}
\label{sec:method}

\subsection{Design principles}
\label{sec:principles}

We state five constraints implied by the aliasing formulation and use them to
motivate each design choice:
\begin{description}
    \item[P1. Perception-level intervention.] Aliasing is a perception-level
    failure; the fix should live in the same layer where the problem occurs.
    \item[P2. Feature-level rewrite.] The disambiguation should deliver a
    corrected current-frame representation, not an additional reasoning stream.
    \item[P3. Sparse sampling.] The task-relevant time constant is seconds; we
    sample $M{=}5$ frames at interval $I{=}10$.
    \item[P4. Per-token gating.] Aliasing is local; fusion must be selective.
    \item[P5. Absolute-time encoding.] Different task phases may share visual
    motifs; an absolute timestamp prevents the module from conflating them.
\end{description}

\subsection{Architecture}
\label{sec:arch}

Let $V \in \mathbb{R}^{B \times L \times V \times N \times D}$ denote the
current-frame DeepStack features, with $B$ batch, $L{=}3$ DeepStack levels,
$V{=}2$ camera views, $N$ patch-merged tokens per view, $D$ feature dimension.
Let $M \in \mathbb{R}^{B \times L \times T \times V \times N \times D}$ denote
the memory tensor with $T$ past timesteps. For each DeepStack level
$\ell \in \{1,\dots,L\}$:

\begin{align}
\widetilde{V}_\ell &= V_\ell + \mathrm{PE}(t),                              \label{eq:pe}\\
H_\ell            &= \mathrm{Attn}_\ell^{\mathrm{read}}\!\big(M_{\ell,-1},\;
                      \mathrm{flatten}_T(M_\ell)\big),                      \label{eq:read}\\
g_\ell            &= \sigma\!\big(W_\ell [H_\ell ;\; \widetilde{V}_\ell]\big), \label{eq:gate}\\
V'_\ell           &= \mathrm{RMSNorm}_\ell\!\big(g_\ell \odot \widetilde{V}_\ell
                      + (1-g_\ell) \odot H_\ell\big).                        \label{eq:fuse}
\end{align}

Eq.~\eqref{eq:pe} injects the absolute timestep, realizing P5.
Eq.~\eqref{eq:read} performs a slot-attention read using the most recent past
frame $M_{\ell,-1}$ as query and the flattened memory tensor as key-value,
realizing P3. Eq.~\eqref{eq:gate} computes a per-token, per-view, per-level
gate, realizing P4. Eq.~\eqref{eq:fuse} produces the fused feature that replaces
$V_\ell$ in the DeepStack tuple handed to the language model, realizing P1--P2.
The parameters $\{W_\ell, \mathrm{Attn}_\ell^{\mathrm{read}}, \mathrm{RMSNorm}_\ell\}$
are independent across levels.

\subsection{Memory construction}
\label{sec:memconstr}

\paragraph{Training.} For a sample at trajectory step $b$ we collect
\[
\mathrm{idx}_i = \max(0,\; b - i \cdot I),\qquad i \in \{M, M{-}1, \dots, 1\}.
\]
This yields an oldest-first window that \emph{excludes} the current frame~$b$
and falls back to the earliest frame when the trajectory is too short.

\paragraph{Rollout.} At inference the \emph{same} rule is applied against the
online history buffer: the current frame is not part of the memory, and padding
repeats the earliest cached frame. See Alg.~\ref{alg:mem} for pseudocode.

\begin{algorithm}[h]
\caption{Online memory construction at rollout step $b$}
\label{alg:mem}
\begin{algorithmic}[1]
\State \textbf{Input:} history buffer $H$ (each entry a list of views), interval $I$, memory length $M$
\State \textbf{Output:} memory frames $\mathcal{M}$ (oldest first, current frame excluded)
\State $b \gets |H| - 1$
\For{$i = M, M-1, \dots, 1$}
    \State $j \gets \max(0,\; b - i \cdot I)$
    \State append $H[j]$ to $\mathcal{M}$
\EndFor
\State \Return $\mathcal{M}$
\end{algorithmic}
\end{algorithm}

\subsection{Integration with the VLM}

The memory forward pass is a single conditional block inside the Qwen3-VL model
forward. Memory images are encoded by the existing vision tower and DeepStack
projector into features of the same shape as $V$; these features are rearranged
into the memory tensor $M$ and consumed by Eqs.~\eqref{eq:pe}--\eqref{eq:fuse};
the resulting $V'$ replaces the original DeepStack tuple passed to the language
model.

\subsection{Action head and training objective}
We use an OFT-style L1 regression head: the hidden states at $C$ action
placeholder tokens are gathered from the last transformer layer, passed through
an MLP, and supervised against normalized expert actions. For STMB the head is
unchanged from the memoryless baseline.

%=========================================================================
\section{Diagnostics: Four Silent Failure Modes}
\label{sec:diagnostics}

Memory-augmented VLA pipelines have a large surface area, and it is easy for
the memory pathway to be miscalibrated without any loud error signal. We
enumerate four failure modes we encountered and explain how STMB avoids each.

\paragraph{F1. Training--evaluation frame-index mismatch.}
The training dataloader collects memory at indices $\{b - M I, \dots, b - I\}$,
which \emph{excludes} the current frame. An unaligned rollout that instead
appends the current frame at the end of the memory shifts every slot by one
interval and changes the most-recent-memory slot from $b - I$ to $b$.
Because the slot-attention write query is $M_{\ell,-1}$, this shift directly
attacks the query--key relationship that the module was trained to exploit.

\paragraph{F2. Zero-image padding.}
When fewer than $M$ past frames are available, a na\"ive rollout pads with an
all-zero RGB image. After RGB normalization such an image is a large-norm
out-of-distribution input that the memory pathway never saw at training, where
padding was implemented by repeating the earliest real frame.

\paragraph{F3. Write-query self-reference.}
If F1 is present (current frame included at $M_{\ell,-1}$), the slot-attention
read degenerates into attending the current frame against a tensor that
includes itself, biasing the gate toward an identity update and effectively
disabling memory.

\paragraph{F4. Causal confusion by history.}
The classical pitfall~\cite{dehaan2019causal}: history correlates with past
actions and is an easy shortcut for the policy. We mitigate this structurally
by (a) sparse sampling at $I{=}10$, which strips fine-grained action traces,
and (b) per-token gating, which keeps the memory contribution local to regions
that actually need disambiguation.

\paragraph{A minimal patch.}
Algorithm~\ref{alg:mem} is the minimal change needed to eliminate F1--F3
simultaneously: it follows the training convention exactly. We report the
effect of this patch as a stand-alone ablation in \S\ref{sec:exp-ablation}.

%=========================================================================
\section{Experiments}
\label{sec:exp}

\subsection{Setup}

\paragraph{Backbone.} Qwen3-VL (RynnBrain-CoP-8B~\cite{alibaba2025rynnbrain}),
loaded with \texttt{bf16} and \texttt{flash\_attention\_2}. The DeepStack
projector is kept active; STMB operates on its output.

\paragraph{Action head.} \texttt{L1RegressionActionHead} with
\texttt{action\_dim}=7, \texttt{state\_dim}=7, chunk $C{=}8$
(future window $7$, past window $0$).

\paragraph{Memory hyperparameters.} $M{=}5$, $I{=}10$, $L{=}3$, $V{=}2$,
hidden dim matching the backbone. Memory bank parameters initialised as in
\S\ref{sec:arch}.

\paragraph{Data.} LIBERO converted to LeRobot v2 / v3 format; combined
\texttt{libero\_all} split for training. Evaluation on the five suites
\textsc{spatial}, \textsc{object}, \textsc{goal}, \textsc{10}, \textsc{90}
with 50 trials per task.

\paragraph{Optimisation.} Per-module learning rates:
\texttt{qwen\_vl\_interface}=$1\mathrm{e}{-}5$,
\texttt{action\_model}=$1\mathrm{e}{-}4$;
cosine schedule with warmup; AdamW ($\beta{=}(0.9, 0.95)$). Accelerate with
DeepSpeed ZeRO-2, gradient checkpointing, mixed precision.

\paragraph{Compute.} \todo{GPU hours / hardware}. \paragraph{Seeds.}
\todo{number of seeds}.

\subsection{Main results on LIBERO}

\begin{table}[h]
\centering
\caption{Success rate (\%) on LIBERO. All rows share the same backbone, action
head, dataloader, and schedule. ``Unaligned'' denotes STMB \emph{without} the
memory-construction patch from \S\ref{sec:diagnostics}; ``Aligned'' is the full
proposed method. Arrow vs.~memoryless baseline in parentheses.}
\label{tab:libero}
\small
\begin{tabular}{lcccccc}
\toprule
Method & \textsc{spatial} & \textsc{object} & \textsc{goal} & \textsc{10} & \textsc{90} & Avg. \\
\midrule
Memoryless baseline (RynnBrain-OFT)    & \num & \num & \num & \num & \num & \num \\
\quad + STMB (unaligned)               & \num & \num & \num & \num & \num & \num \\
\quad + STMB (aligned, \textbf{ours})  & \num & \num & \num & \num & \num & \num \\
\midrule
MemoryVLA~\cite{shi2025memoryvla}      & \num & \num & \num & \num & \num & \num \\
PAM~\cite{pam2025}                     & \num & \num & \num & \num & \num & \num \\
\bottomrule
\end{tabular}
\end{table}

\subsection{Aliasing-stratified evaluation}
\label{sec:exp-alias}

We bucket the 130 LIBERO tasks by task-average aliasing score
(Def.~\ref{def:alias}, computed with a frozen Qwen3-VL DeepStack encoder, $k{=}32$,
action threshold $\delta$ set at the 75\textsuperscript{th} percentile of
inter-action distance). We then report gains of STMB over the memoryless
baseline stratified into low / medium / high aliasing buckets.

\begin{table}[h]
\centering
\caption{Aliasing-stratified gains. We expect low aliasing: parity; high
aliasing: largest gains.}
\label{tab:alias-bucket}
\small
\begin{tabular}{lcccc}
\toprule
Aliasing bucket & \# tasks & Baseline SR & + STMB SR & $\Delta$ \\
\midrule
Low    & \num & \num & \num & \num \\
Medium & \num & \num & \num & \num \\
High   & \num & \num & \num & \num \\
\bottomrule
\end{tabular}
\end{table}

\begin{figure}[h]
\centering
\todo{Scatter plot: x-axis = task aliasing score, y-axis = STMB gain over
baseline, one point per LIBERO task. Overlay a linear fit and report the
correlation coefficient.}
\caption{Per-task correlation between aliasing score and STMB gain.}
\label{fig:alias-scatter}
\end{figure}

\subsection{Ablations}
\label{sec:exp-ablation}

We ablate each of P1--P5 and the memory-construction patch independently.

\begin{table}[h]
\centering
\caption{Ablations on LIBERO (average success rate \%).}
\label{tab:ablations}
\small
\begin{tabular}{lc}
\toprule
Variant & Avg.~SR \\
\midrule
STMB full (ours)                                                 & \num \\
(P3) $M{=}1$ (single past frame)                                  & \num \\
(P3) $M{=}3$                                                      & \num \\
(P3) $M{=}7$                                                      & \num \\
(P3) Interval $I{=}5$                                              & \num \\
(P3) Interval $I{=}20$                                             & \num \\
(P4) Remove gate (plain residual fusion $V' = V + H$)             & \num \\
(P5) Remove temporal PE                                           & \num \\
Single-view memory (agentview only)                               & \num \\
Single-view memory (wrist only)                                   & \num \\
Unaligned memory construction (reproduces F1+F2+F3)               & \num \\
\bottomrule
\end{tabular}
\end{table}

\subsection{Failure-mode study (F1--F4)}

We construct targeted diagnostic experiments: (i)~force the memory at eval to
include the current frame (F1/F3), (ii)~replace the earliest-frame padding with
zeros (F2), (iii)~subsample at $I{=}1$ to stress-test causal confusion (F4).
\todo{Results table and discussion.}

%=========================================================================
\section{Discussion}
\label{sec:discussion}

\paragraph{Why perception-level rather than reasoning-level memory?}
Aliasing is a property of observations, not of task plans. Fixing it at the
language-model input (with extra memory tokens) pays a sequence-length cost
to solve a problem that sits one layer lower. STMB writes corrections directly
into the pathway where the problem manifests.

\paragraph{Relation to long-horizon VLAs.}
STMB is deliberately short-term. Tasks whose causal support exceeds $M \cdot I$
frames are out of scope; they are better served by architectures such as
MemoryVLA and LoLA, which STMB is complementary to (see \S\ref{sec:related}).

\paragraph{Limitations.}
(i)~STMB is currently specialised to DeepStack-style VLMs; extending to
single-projector backbones is future work. (ii)~The aliasing score uses a
fixed visual encoder; a learned aliasing score could be more predictive.
(iii)~We do not address non-visual aliasing (e.g., identical visuals but
different proprioception); joint state--image memory is a natural extension.

%=========================================================================
\section{Conclusion}

We named and formalised \emph{visual observation aliasing}, a silent but
pervasive failure mode of Markov VLA policies. Reading the problem back off its
definition gave us five concrete design constraints, which we instantiated in
STMB, a feature-level memory adapter that writes directly into the DeepStack
visual pathway of a VLM backbone. A diagnostic catalogue of four silent
failure modes, together with a one-algorithm alignment patch, makes the
pipeline reproducible.
Aliasing-stratified evaluation on LIBERO \todo{ties the improvement to a
property of the data}, suggesting that the right question for
memory-augmented VLAs is not ``how much history'' but ``what kind of
ambiguity does the history dissolve.''

%=========================================================================
\bibliographystyle{plain}
\begin{thebibliography}{99}

\bibitem{brohan2023rt2}
A.~Brohan et al.
\newblock RT-2: Vision-Language-Action Models Transfer Web Knowledge to Robotic Control.
\newblock \emph{CoRL}, 2023.

\bibitem{kim2024openvla}
M.~Kim et al.
\newblock OpenVLA: An Open-Source Vision-Language-Action Model.
\newblock \emph{CoRL}, 2024.

\bibitem{black2024pi0}
K.~Black et al.
\newblock $\pi_0$: A Vision-Language-Action Flow Model for General Robot Control.
\newblock arXiv:2410.24164, 2024.

\bibitem{bai2023qwen}
J.~Bai et al.
\newblock Qwen-VL: A Versatile Vision-Language Model for Understanding, Localization, Text Reading, and Beyond.
\newblock arXiv:2308.12966, 2023.

\bibitem{locatello2020slot}
F.~Locatello et al.
\newblock Object-Centric Learning with Slot Attention.
\newblock \emph{NeurIPS}, 2020.

\bibitem{dehaan2019causal}
P.~de~Haan, D.~Jayaraman, S.~Levine.
\newblock Causal Confusion in Imitation Learning.
\newblock \emph{NeurIPS}, 2019.

\bibitem{causalbot2025}
Anonymous.
\newblock Improving Generalization Ability of Robotic Imitation Learning by Resolving Causal Confusion in Observations.
\newblock arXiv:2507.22380, 2025.

\bibitem{shi2025memoryvla}
S.~Shi et al.
\newblock MemoryVLA: Perceptual-Cognitive Memory in Vision-Language-Action Models for Robotic Manipulation.
\newblock arXiv:2508.19236, 2025.

\bibitem{remem2025}
Anonymous.
\newblock Empowering Vision-Language-Action Model with Memory via Dual-Level Recurrent Queries.
\newblock arXiv:2603.12942, 2025.

\bibitem{pam2025}
Anonymous.
\newblock Resolving State Ambiguity in Robot Manipulation via Adaptive Working Memory Recoding.
\newblock arXiv:2512.24638, 2025.

\bibitem{lola2025}
Anonymous.
\newblock LoLA: Long Horizon Latent Action Learning for General Robot Manipulation.
\newblock arXiv:2512.20166, 2025.

\bibitem{zhao2025cotvla}
Y.~Zhao et al.
\newblock CoT-VLA: Visual Chain-of-Thought Reasoning for Vision-Language-Action Models.
\newblock \emph{CVPR}, 2025.

\bibitem{alibaba2025rynnbrain}
Alibaba DAMO Academy.
\newblock RynnBrain: A Vision-Language Model Family for Embodied Agents.
\newblock Technical Report, 2025.

\end{thebibliography}

\end{document}
